\begin{tikzpicture}[
    x=1cm,
    y=1cm,
    entrycell/.style={
        draw,
        semithick,
        align=left,
        text width=31mm,
        minimum height=5.8mm,
        inner sep=2pt,
        font=\scriptsize
    },
    codecell/.style={
        draw,
        semithick,
        align=left,
        text width=39mm,
        minimum height=5.8mm,
        inner sep=2pt,
        font=\scriptsize
    },
    entryfunc/.style={entrycell, fill=orange!18},
    entryoperand/.style={entrycell, fill=orange!6},
    codefunc/.style={codecell, fill=orange!18},
    codeoperand/.style={codecell, fill=orange!6},
    markerentry/.style={entrycell, fill=gray!8, align=center},
    markercode/.style={codecell, fill=gray!8, align=center},
    handler/.style={
        draw,
        rounded corners=2pt,
        semithick,
        align=left,
        text width=47mm,
        minimum height=8mm,
        inner sep=4pt,
        fill=white,
        font=\scriptsize
    },
    ipbox/.style={
        draw,
        rounded corners=2pt,
        semithick,
        align=center,
        minimum width=8mm,
        minimum height=5mm,
        inner sep=4pt,
        fill=gray!8,
        font=\scriptsize
    },
    flow/.style={->, semithick},
    targetflow/.style={->, semithick, rounded corners=4pt},
    returnflow/.style={->, semithick, densely dashed}
]

\node[font=\small\bfseries, align=center] at (-4.3,0.85) {\\\texttt{entry\_points}};
\node[font=\small\bfseries, align=center] at (0.75,0.85) {\\\texttt{all\_code}};
\node[font=\small\bfseries] at (8.1,0.65) {\texttt{ops.c}};

\node[entryfunc] (entrycall1) at (-4.3,0) {\texttt{func: op\_call}};
\node[entryoperand, below=0pt of entrycall1] (entrytarget1) {\texttt{target: main\_1}};
\node[entryoperand, below=0pt of entrytarget1] (entryargs1) {\texttt{num: 0}};
\node[entryfunc,    below=0pt of entryargs1] (entrydrop1) {\texttt{func: op\_drop}};
\node[markerentry,  below=0pt of entrydrop1] (entrygap) {$\cdots$};
\node[entryfunc,    below=0pt of entrygap]   (entrycalln) {\texttt{func: op\_call}};
\node[entryoperand, below=0pt of entrycalln] (entrytargetn) {\texttt{target: main\_n}};
\node[entryoperand, below=0pt of entrytargetn] (entryargsn) {\texttt{num: 0}};
\node[entryfunc,    below=0pt of entryargsn] (entrydropn) {\texttt{func: op\_drop}};
\node[entryfunc,    below=0pt of entrydropn] (entryeof) {\texttt{func: op\_eof}};


\node[codefunc] (mainbegin) at (0.75,0) {\texttt{func: op\_begin (main\_1)}};
\node[codefunc,    below=0pt of mainbegin]  (constop) {\texttt{func: op\_const}};
\node[codeoperand, below=0pt of constop]    (constarg) {\texttt{num: 42}};
\node[codefunc,    below=0pt of constarg]   (callop) {\texttt{func: op\_call}};
\node[codeoperand, below=0pt of callop]     (calltarget) {\texttt{target: f}};
\node[codeoperand, below=0pt of calltarget] (callargs) {\texttt{num: n}};
\node[codefunc,    below=0pt of callargs]   (aftercall) {следующий обработчик};
\node[markercode,  below=0pt of aftercall]  (gap2) {$\cdots$};
\node[codefunc,    below=0pt of gap2]       (fbegin) {\texttt{func: op\_begin (f)}};
\node[markercode,  below=0pt of fbegin]     (fbody) {$\cdots$ тело функции};
\node[codefunc,    below=0pt of fbody]      (endop) {\texttt{func: op\_end}};

\node[ipbox] (ip) at ($(callop.west)+(-0.9,0)$) {\texttt{ip}};

\node[handler] (hcall) at (8.1,-1.25)
    {\texttt{void op\_call(DECL\_STATE) \{}\\
     \texttt{\hspace*{1.2em}ip++;}\\
     \texttt{\hspace*{1.2em}target = ip->target; ip++;}\\
     \texttt{\hspace*{1.2em}n\_args = ip->num; ip++;}\\
     \texttt{\hspace*{1.2em}PUSH\_FRAME(..., ip,}\\
     \texttt{\hspace*{2.4em}sp + n\_args);}\\
     \texttt{\hspace*{1.2em}ip = target;}\\
     \texttt{\hspace*{1.2em}DISPATCH\_JUMP();}\\
     \texttt{\}}};

\node[font=\large] (handlergap) at (8.1,-3.45) {$\cdots$};

\node[handler] (hconst) at (8.1,-4.35)
    {\texttt{void op\_const(DECL\_STATE) \{...\}}};

\node[handler] (hbegin) at (8.1,-5.50)
    {\texttt{void op\_begin(DECL\_STATE) \{...\}}};

\node[handler] (hend) at (8.1,-6.65)
    {\texttt{void op\_end(DECL\_STATE) \{...\}}};

\draw[flow] (ip.east) -- (callop.west);
\draw[flow] (callop.east) -- (hcall.west);
\draw[flow] (mainbegin.east) -- (hbegin.west);
\draw[flow] (constop.east) -- (hconst.west);
\draw[flow] (endop.east) -- (hend.west);

\draw[targetflow] (entrytarget1.east) -- (mainbegin.west);
\draw[targetflow] (entrygap.east) -- (gap2.west);

\draw[targetflow] (entrytargetn.east) -- (gap2.west);
\draw[targetflow] (calltarget.east) -- ++(0.55,0) |- (fbegin.east);
\draw[returnflow] (hend.west) to[out=180,in=0] (aftercall.east);

\end{tikzpicture}
